\section{Erd\H{o}s Problem 45: a divisor construction and its growth scale}

\subsection{1) Problem and known input}
For every $k\geq2$, does there exist an integer $n_k$ such that every
$k$-colouring of
\[
 D(n_k)=\{d:1<d<n_k,\ d\mid n_k\}
\]
contains a monochromatic subset $S$ with $\sum_{d\in S}1/d=1$?
The current problem page, checked on 23 September 2026, records an
affirmative answer from Croot's unit-fraction colouring theorem.

The precise external input is the following finite statement from
Croot's 2003 paper: there is an absolute constant $b>1$ such that
every $k$-colouring of the integers in $[2,b^k]$ has a monochromatic
set of distinct denominators whose reciprocals sum to $1$.
The analytic proof of that theorem is not reproduced here. We give
the full transfer to proper divisors and a matching growth scale,
with each further input identified. This is a reconstruction of known
results, not a claim of a new solution.

\subsection{2) Explicit transfer from an interval to proper divisors}
Choose an integer $L\geq\max(3,b^k)$ and put
\[
 n_k=\operatorname{lcm}(1,2,\ldots,L).
\]
Every integer from $2$ through $L$ divides $n_k$. These divisors really
are proper: since $L$ and $L-1$ are coprime,
\[
 n_k\geq L(L-1)>L.
\]
Restrict any given colouring of $D(n_k)$ to $\{2,\ldots,L\}$.
Croot's theorem supplies a monochromatic $S\subseteq\{2,\ldots,L\}$
with reciprocal sum $1$. Because this entire interval lies in
$D(n_k)$, this is exactly the subset required in the problem.
No repeated denominator or the excluded divisor $1$ has been used.

Even the elementary estimate $n_k\leq L!\leq L^L$ gives a doubly
exponential bound. With $L=\lceil\max(3,b^k)\rceil$,
\[
 \log n_k\leq L\log L\leq C^k
\]
for a suitable absolute $C>1$ and all $k\geq2$, enlarging $C$ to
cover the finitely many initial cases. Thus $n_k\leq\exp(C^k)$.
One can sharpen the estimate on this particular least common multiple,
but that is unnecessary for the stated growth scale.

\subsection{3) A rigorous greedy obstruction for small reciprocal mass}
Here is a useful sufficient condition for a colouring with no solution.
If a finite set of denominators all at least $2$ has total reciprocal
mass less than $k/2$, assign its elements successively to $k$ colour
classes, keeping the sum in each class strictly below $1$.

To justify that the assignment never gets stuck, let the next weight
be $w=1/d\leq1/2$. If every class rejected it, each existing class
would have mass at least $1-w\geq1/2$. The already assigned mass
would then be at least $k/2$, contradicting the assumed total mass.
Every class consequently has mass less than $1$, and no subset of
one class can have mass exactly $1$.

Applying this to $D(n)$ shows that any integer having the required
Ramsey property must satisfy
\[
 \sum_{d\in D(n)}\frac1d\geq\frac k2.
\]
This deliberately coarse necessary bound is enough below. We do not
need a sharper packing assertion with threshold $k$.

\subsection{4) A doubly exponential necessary size}
We use the classical Mertens product estimate
\[
 \prod_{p\leq y}(1-1/p)^{-1}\ll\log y\qquad(y\geq3)
\]
as a second external theorem. For $n$ sufficiently large, put
$y=\log n$. Factorization of the divisor sum gives
\[
 \sum_{d\mid n}\frac1d
 =\prod_{p^a\parallel n}(1+p^{-1}+\cdots+p^{-a})
 \leq\prod_{p\mid n}(1-1/p)^{-1}.
\]
The contribution of $p\leq y$ is $O(\log y)$ by Mertens.
If there are $r$ distinct prime factors exceeding $y$, then
$y^r<n$, so $r\leq\log n/\log y$. Since
$-\log(1-1/p)\leq2/p$,
\[
 \log\prod_{\substack{p\mid n\\p>y}}(1-1/p)^{-1}
 \leq\frac{2r}{y}
 \leq\frac{2}{\log\log n}.
\]
The large-prime contribution is bounded by an absolute constant.
Consequently
\[
 \sum_{d\mid n}\frac1d\leq C_0\log\log n
\]
for all sufficiently large $n$, with an absolute $C_0$.
Combining this with the necessary reciprocal mass gives
\[
 \log\log n_k\geq\frac{k}{2C_0},\qquad
 n_k\geq\exp\bigl(\exp(c k)\bigr)
\]
for an absolute $c>0$ and sufficiently large $k$.
Together with the construction, this identifies the doubly exponential
scale. It does not identify the least possible integer $n_k$ or the
optimal constants in the exponent.

\subsection{5) Verification and outcome}
The proper-divisor check $n_k>L$ prevents an endpoint error in the
interval transfer. The lower bound uses only positive weights, so a
class of mass less than $1$ cannot contain a hidden exact solution.
Both uses of asymptotic constants are uniform in $k$; Croot's $b$ and
Mertens' constant are absolute.

\textbf{Outcome: SOLVED using Croot's finite colouring theorem.}
The Ramsey transfer is complete. The additional lower bound uses the
explicitly stated Mertens estimate. Neither analytic theorem is
asserted to have been independently reproved or formally checked here.

\subsection{6) Sources}
\begin{itemize}
\item E. S. Croot III, \emph{On a coloring conjecture about unit
fractions}, Annals of Mathematics 157 (2003), 545--556, introductory
corollary, \texttt{https://annals.math.princeton.edu/wp-content/uploads/annals-v157-n2-p04.pdf}.
\item T. F. Bloom, Problem 45,
\texttt{https://www.erdosproblems.com/45}, checked 23 September 2026.
The page credits the lower-bound observation to Mehtaab Sawhney;
the argument above uses a weaker elementary packing threshold which
already gives the same scale.
\item Mertens' product estimate is also stated as equation (2) in
T. F. Bloom, \emph{On a density conjecture about unit fractions},
\texttt{https://arxiv.org/html/2112.03726v2}.
\end{itemize}

The estimate records full coverage with the named external theorems
accepted; it is not a novelty claim.
\subsection{7) COMPLETION ESTIMATE}
\textbf{COMPLETION: 100\%}
